\documentclass[10pt]{article}
\usepackage[margin=0.75in]{geometry}
\usepackage{amsmath,amssymb,amsfonts}
\usepackage{booktabs}
\usepackage{graphicx}
\usepackage{enumitem}
\setlist{nosep,leftmargin=*}
\usepackage{hyperref}
\usepackage{tikz}
\usepackage{caption}
\usepackage{subcaption}
\usepackage{xcolor}
\usepackage{float}
\usepackage{titlesec}
\titlespacing*{\section}{0pt}{1.2ex plus 0.3ex}{0.5ex plus 0.1ex}
\titlespacing*{\subsection}{0pt}{0.8ex plus 0.2ex}{0.3ex plus 0.1ex}
\setlength{\parskip}{0.2ex plus 0.1ex}
\setlength{\parindent}{1em}
\setlength{\abovecaptionskip}{3pt}
\setlength{\belowcaptionskip}{-4pt}
\setlength{\textfloatsep}{6pt}
\setlength{\floatsep}{6pt}
\setlength{\intextsep}{6pt}

\title{\textbf{Neural SC Decoder for Gaussian MAC Polar Codes:\\Problem Statement and Request for Ideas}}
\author{Project Documentation for External Collaboration}
\date{March 2026}

\begin{document}
\maketitle

%% ============================================================
%% PAGE 1: BACKGROUND
%% ============================================================
\section{Background}

\subsection{Two-User MAC Polar Codes}

Polar codes, introduced by Ar{\i}kan~\cite{arikan2009}, achieve channel capacity through a phenomenon called \emph{polarization}: as block length $N = 2^n$ grows, each synthetic sub-channel becomes either perfectly reliable or perfectly noisy. For a two-user multiple-access channel (MAC), \"Onay~\cite{onay2013} showed that polar codes can simultaneously achieve the capacity region by having each user independently encode with a polar code, while a joint successive cancellation (SC) decoder processes the superimposed codeword.

\subsection{Channel Model: Gaussian MAC}

We consider a two-user Gaussian MAC where users $X$ and $Y$ transmit BPSK-modulated polar codewords. The received signal is:
\begin{equation}
Z = (1 - 2X) + (1 - 2Y) + W, \qquad W \sim \mathcal{N}(0, \sigma^2),
\label{eq:gmac}
\end{equation}
where $X, Y \in \{0, 1\}$ are coded bits and the per-user SNR is $\text{SNR} = 1/\sigma^2$. The observation $Z \in \mathbb{R}$ is continuous, unlike the binary-erasure MAC (BEMAC) where $Z = X + Y \in \{0, 1, 2\}$.

At each leaf position $t$ of the decoding tree, the decoder must compute the posterior over the four joint hypotheses $(u_t, v_t) \in \{00, 01, 10, 11\}$, where $u_t$ and $v_t$ are the $t$-th information bits of users $X$ and $Y$ respectively.

\subsection{Decoding Paths and Code Classes}

The joint SC decoder processes $2N$ bits total ($N$ per user) in a specific interleaving order called a \emph{path}. Let $\pi_i$ index the path, with $0 \le \pi_i \le N$. Three classes arise:

\begin{itemize}[nosep]
    \item \textbf{Class A} ($\pi_i = 0$): Decode all of $Y$ first, then all of $X$. An \emph{extreme} path.
    \item \textbf{Class C} ($\pi_i = N$): Decode all of $X$ first, then all of $Y$. Also extreme.
    \item \textbf{Class B} ($\pi_i = N/2$): Interleave $X$ and $Y$ decisions. A \emph{monotone chain} path.
\end{itemize}

Classes A and C are symmetric and use only \texttt{CalcLeft} and \texttt{CalcRight} (downward tree operations). Class B additionally requires \texttt{CalcParent}---an \emph{upward} operation that propagates decided information back up the tree before descending into the next user's sub-tree. This is the key algorithmic difference.

\subsection{Complexity and Motivation}

The analytical SC decoder for the two-user MAC has complexity:
\begin{equation}
\mathcal{O}(m^2 \cdot N \log N),
\end{equation}
where $m = 4$ is the joint alphabet size $|\{00, 01, 10, 11\}|$. Each tree operation manipulates probability tensors of size $m^{\ell}$ at depth $\ell$, but the $\mathcal{O}(N \log N)$ computational graph structure keeps the total cost manageable.

For channels \emph{with memory} (e.g., inter-symbol interference, fading), each tree operation must track a state space of size $S$, giving complexity:
\begin{equation}
\mathcal{O}(S^3 \cdot N \log N).
\end{equation}
Even moderate state spaces ($S = 16$) make this prohibitive. Our goal is to replace the $\mathcal{O}(S^3)$ tensor operations with fixed-size neural networks of complexity $\mathcal{O}(md)$, where $d$ is the embedding dimension:
\begin{equation}
\mathcal{O}(S^3 \cdot N \log N) \;\longrightarrow\; \mathcal{O}(md \cdot N \log N).
\end{equation}
This would yield a \textbf{channel-independent decoder}: the tree operations are learned once and reused across any channel, with only the front-end embedding layer changing per channel.

%% ============================================================
%% PAGE 2: ARCHITECTURE
%% ============================================================
\section{Architecture: Neural Computational Graph (NCG) Decoder}

\subsection{Overview}

The NCG decoder replaces each analytical tree operation with a small MLP, while preserving the exact SC decoding schedule (the order in which tree nodes are visited). Every edge in the binary tree stores a $d$-dimensional embedding vector instead of a probability tensor. All MLPs share weights across tree levels, making the decoder \textbf{independent of block length $N$}.

\subsection{Tree Structure}

The binary tree has $N - 1$ internal vertices and $N$ leaves. Edges are indexed $1, \ldots, 2N-1$:
\begin{itemize}[nosep]
    \item Edge 1 = root (channel observation embedding).
    \item Edges $N, \ldots, 2N-1$ = leaves (one per bit position).
    \item Vertex $\beta$ has parent edge $\beta$, left child edge $2\beta$, right child edge $2\beta + 1$.
\end{itemize}

\begin{figure}[H]
\centering
\begin{small}
\begin{verbatim}
             [Edge 1: root]           <-- z_encoder output (d=16)
            /              \
      [Edge 2]           [Edge 3]     <-- CalcLeft / CalcRight (3d -> d)
      /      \           /      \         CalcParent goes UP (2d -> d)
  [Edge 4] [Edge 5] [Edge 6] [Edge 7]
     |        |        |        |
   Leaf 0  Leaf 1   Leaf 2   Leaf 3   <-- Emb2Logits (d -> 4 classes)
\end{verbatim}
\end{small}
\vspace{-6pt}
\caption{NCG tree for $N = 4$. Each edge stores a $d$-dimensional embedding.}
\label{fig:tree}
\end{figure}

\subsection{Neural Operations}

\begin{enumerate}[nosep]
    \item \textbf{CalcLeft} (top-down, $\mathbb{R}^{3d} \to \mathbb{R}^d$): Given parent embedding $\mathbf{e}_\beta$ and child embeddings, compute left child's top-down message. Input is concatenation $[\mathbf{e}_\beta \| \mathbf{e}_{2\beta} \| \mathbf{e}_{2\beta+1}]$. Two-layer MLP with ReLU.

    \item \textbf{CalcRight} (top-down, $\mathbb{R}^{3d} \to \mathbb{R}^d$): Same structure as CalcLeft but separate weights. Computes right child's top-down message.

    \item \textbf{CalcParent} (bottom-up, $\mathbb{R}^{2d} \to \mathbb{R}^d$): Propagates a decision back \emph{up} the tree. Takes child embedding (after decision) and sibling embedding, produces updated parent embedding. Uses a \textbf{gated residual MLP}: $\mathbf{e}_\text{out} = \alpha \cdot \text{MLP}([\mathbf{e}_\text{child} \| \mathbf{e}_\text{sibling}]) + (1 - \alpha) \cdot \mathbf{e}_\text{parent}$, where $\alpha = \sigma(\text{gate}(\cdot))$.

    \item \textbf{Emb2Logits} ($\mathbb{R}^d \to \mathbb{R}^4$): At each leaf, converts embedding to 4-class log-probabilities over $(u, v) \in \{00, 01, 10, 11\}$.

    \item \textbf{Logits2Emb} ($\mathbb{R}^4 \to \mathbb{R}^d$): After a leaf decision, converts the decided log-probabilities back to an embedding for tree propagation.
\end{enumerate}

\subsection{Channel Front-Ends}

The front-end converts channel observations into root embeddings:
\begin{itemize}[nosep]
    \item \textbf{BEMAC}: \texttt{nn.Embedding(3, d=16)} --- discrete lookup for $Z \in \{0, 1, 2\}$.
    \item \textbf{GMAC}: \texttt{MLP}$(1 \to 32 \to 16)$ --- continuous mapping for $Z \in \mathbb{R}$.
\end{itemize}
The GMAC front-end must learn to map a continuous Gaussian observation to the same embedding space that the tree operations expect. This is a key difficulty.

\subsection{Training}

\begin{itemize}[nosep]
    \item \textbf{Teacher forcing}: At training time, leaf decisions use ground-truth $(u_t, v_t)$ to avoid error propagation. Loss is cross-entropy on the 4-class output at each leaf.
    \item \textbf{Curriculum learning}: Essential for convergence. Train at $N = 8$ from scratch, then fine-tune at $N = 16 \to 32 \to 64 \to \cdots \to 1024$.
    \item \textbf{Total parameters}: $\sim$39K (all shared across tree levels).
    \item \textbf{Training cost}: $\sim$1 second/iteration at $N = 128$ on CPU; $\sim$4 seconds at $N = 1024$.
\end{itemize}

\subsection{Soft-Bit Bridge (for CalcParent)}

For the BEMAC channel, we found that a pure neural CalcParent fails. Instead, we use an analytical \emph{bridge}: convert embedding to logits (Emb2Logits), apply analytical circular convolution in probability space (differentiable), then convert back (Logits2Emb). This preserves exact information flow while keeping end-to-end differentiability.

%% ============================================================
%% PAGE 3: RESULTS
%% ============================================================
\section{Results: What Works and What Doesn't}

\subsection{BEMAC Results (Discrete Channel, $Z = X + Y$)}

The NCG decoder with Soft-Bit Bridge for CalcParent achieves remarkable results on the discrete BEMAC:

\begin{table}[H]
\centering
\caption{BEMAC Class B: NCG vs.\ Analytical SC (SNR matched, proper frozen sets)}
\label{tab:bemac}
\begin{tabular}{@{}rccl@{}}
\toprule
$N$ & NN BLER & SC BLER & Ratio \\
\midrule
16  & 0.011  & 0.011  & 1.0$\times$ (matches) \\
32  & 0.009  & 0.008  & 1.1$\times$ \\
64  & 0.003  & 0.006  & 0.5$\times$ \textbf{(NN wins)} \\
128 & 0.001  & 0.002  & 0.6$\times$ \textbf{(NN wins)} \\
512 & 0.0004 & 0.0004 & 1.0$\times$ (matches) \\
1024 & 0.0003 & 0.0005 & 0.6$\times$ \textbf{(NN wins)} \\
\bottomrule
\end{tabular}
\end{table}

The NN \textbf{matches or beats SC at every tested block length}. The architecture and training methodology are sound. The Soft-Bit Bridge for CalcParent works perfectly for discrete channels.

\subsection{Neural SCL: List Decoding on Top of Neural Tree (NEW)}

A major breakthrough: we implemented \textbf{Neural SC-List (SCL) decoding} on top of the neural tree operations. Instead of greedy argmax at each leaf, the decoder maintains $L$ candidate paths through the tree, forking at each non-frozen leaf into up to $2L$ candidates (one per bit hypothesis, marginalising over the other user) and keeping the best $L$ by cumulative log-probability.

\begin{table}[H]
\centering
\caption{BEMAC Class B: NN-SCL ($L = 4$) vs.\ SC and SCL-32}
\label{tab:bemac-scl}
\begin{tabular}{@{}rccccl@{}}
\toprule
$N$ & SC & SCL-32 & NN-SC & NN-SCL ($L\!=\!4$) & Ratio to SC \\
\midrule
32  & 0.0080 & 0.0082 & 0.0088 & 0.0073 & 0.92$\times$ \\
64  & 0.0056 & 0.0010 & 0.0030 & 0.00067 & \textbf{0.12$\times$ (8.4$\times$ better)} \\
128 & 0.0020 & 0.0006 & 0.0012 & 0.00067 & \textbf{0.33$\times$ (3$\times$ better)} \\
\bottomrule
\end{tabular}
\end{table}

Key findings:
\begin{itemize}[nosep]
    \item At $N = 64$, NN-SCL ($L\!=\!4$) achieves BLER $= 6.7 \times 10^{-4}$, which is \textbf{8.4$\times$ better than SC} and \textbf{1.5$\times$ better than SCL-32}.
    \item At $N = 128$, NN-SCL matches SCL-32 performance (both $\sim 6 \times 10^{-4}$), and is \textbf{3$\times$ better than SC}.
    \item The neural tree operations learn a richer representation than the analytical probability tensors, which the list decoder can exploit more effectively.
    \item This is the \textbf{first demonstration of neural list decoding for MAC polar codes} outperforming both greedy neural and classical analytical decoders.
\end{itemize}

\subsection{GMAC Results (Continuous Channel, SNR = 6\,dB)}

The same architecture with a continuous front-end fails dramatically at large $N$:

\begin{table}[H]
\centering
\caption{GMAC Class B: NCG vs.\ Analytical SC (SNR = 6\,dB)}
\label{tab:gmac}
\begin{tabular}{@{}rccl@{}}
\toprule
$N$ & NN BLER & SC BLER & Ratio \\
\midrule
32   & 0.040  & 0.046  & 0.87$\times$ \textbf{(NN wins)} \\
64   & 0.031  & 0.025  & 1.25$\times$ \\
128  & 0.032  & 0.015  & 2.1$\times$ \\
256  & 0.025  & 0.005  & 5.0$\times$ \\
512  & 0.045  & 0.001  & 45$\times$ \\
1024 & 0.069  & 0.0003 & 230$\times$ \\
\bottomrule
\end{tabular}
\end{table}

\textbf{The NN BLER stays flat} at $\sim$0.03--0.07 regardless of $N$, while SC BLER drops exponentially (the hallmark of polarization). The gap grows exponentially: $230\times$ at $N = 1024$.

\subsection{Interpretation}

The contrast is stark. On BEMAC (3-valued discrete input), the NN captures polarization perfectly. On GMAC (continuous input), it does not. The only architectural difference is the front-end: \texttt{nn.Embedding(3, 16)} vs.\ \texttt{MLP(1, 32, 16)}. The tree operations are identical.

This suggests the continuous front-end fails to produce embeddings with sufficient \emph{precision} for the tree operations to exploit polarization at large $N$. Polarization requires that information accumulates multiplicatively through $\log N$ tree levels. Any imprecision in the root embedding is amplified exponentially.

%% ============================================================
%% PAGE 4: WHAT WE'VE TRIED
%% ============================================================
\section{What We've Tried (and Why It Failed)}

\subsection{More Training Time}

Diminishing returns. At $N = 1024$, each iteration takes $\sim$4 seconds (2048 sequential MLP calls through the tree). We managed $\sim$4K iterations. At $N = 32$, we ran 30K iterations. The loss converges but BLER does not improve beyond the plateau.

\subsection{Knowledge Distillation from Analytical SC}

Three variants attempted:

\begin{enumerate}[nosep]
    \item \textbf{Phase-based} (freeze all except CalcParent, train CalcParent from teacher): The z\_encoder never trains because it is frozen. Wasted iterations.
    \item \textbf{MSE on all intermediate embeddings}: Teacher and student learn different embedding spaces. MSE loss exploded to $10^{18}$.
    \item \textbf{Clamped MSE}: Stable training but no improvement over baseline (BLER 0.072 vs.\ 0.070).
\end{enumerate}

\textbf{Why it fails}: The analytical decoder operates on probability tensors in $\mathbb{R}^{4^\ell}$ at depth $\ell$. The neural decoder uses $\mathbb{R}^{d}$ embeddings. There is no natural correspondence between these representation spaces.

\subsection{BEMAC$\to$GMAC Transfer Learning}

Freeze the pre-trained BEMAC tree weights (which work perfectly); train only the GMAC front-end.

\begin{itemize}[nosep]
    \item $N = 32$: Works (BLER = 0.095 in 8K iters).
    \item $N = 64$: Fails (loss stalls at 0.81, near random).
\end{itemize}

\textbf{Why it fails}: The tree operations are not truly channel-independent. At larger depths, the embeddings carry channel-specific structure that the GMAC front-end cannot replicate.

\subsection{Fast Cross-Entropy (Auxiliary Loss at Internal Edges)}

Following the NPD literature~\cite{crisp2023}, we added auxiliary cross-entropy losses at every internal tree edge (not just leaves). This provides denser gradient signal.

\textbf{Result}: Helps at $N = 32$, hurts at $N \ge 64$. Internal edges have ambiguous targets when the tree walk is path-dependent.

\subsection{Architecture Variations}

\begin{itemize}[nosep]
    \item Embedding dimension $d$: tested 8, 16, 32, 48. No significant difference.
    \item Front-end hidden size: tested 16 to 128. No difference.
    \item Per-level weights (separate MLPs per tree depth): Needs $10\times$ more training data, untested at scale.
    \item Temperature scaling ($\tau = 1.0, 1.5, 2.0$): $\tau = 1.0$ is optimal; higher $\tau$ hurts.
\end{itemize}

\subsection{Training Tricks}

\begin{itemize}[nosep]
    \item \textbf{Label smoothing}: Helps from-scratch training but not with curriculum.
    \item \textbf{Scheduled sampling}: Good training loss but BLER = 1.0 at test time. The model becomes dependent on noisy inputs during training and cannot recover.
    \item \textbf{Multi-SNR training}: Not yet tested. Could help generalization.
\end{itemize}

\subsection{Pure Neural CalcParent (No Soft-Bit Bridge)}

Attempted a GRU-gated MLP for CalcParent with knowledge distillation from the analytical bridge:
\begin{itemize}[nosep]
    \item Phase A (teacher available): BLER = 0.012. Works.
    \item Phase B/C (teacher removed): Collapses to BLER = 1.0.
\end{itemize}
The neural CalcParent cannot maintain information fidelity without the analytical bridge.

%% ============================================================
%% PAGE 5: ANALYSIS AND OPEN QUESTIONS
%% ============================================================
\section{Analysis and Open Questions}

\subsection{The Core Problem: Failure to Learn Polarization}

The defining feature of polar codes is \emph{polarization}: as $N \to \infty$, the fraction of ``good'' sub-channels approaches the channel capacity, and each good sub-channel becomes arbitrarily reliable. The analytical SC decoder exploits this: its BLER drops as $\mathcal{O}(2^{-N^{0.5}})$.

Our neural decoder's BLER stays flat. It does not learn to exploit the exponentially increasing reliability of polarized sub-channels. This is not an architecture problem per se---the BEMAC results prove the architecture works---but a \emph{representation problem} specific to continuous channels.

\subsection{Gradient Analysis}

At $N \ge 64$, we observe:
\begin{itemize}[nosep]
    \item \textbf{z\_encoder (front-end)}: Receives near-zero gradients. The 10+ tree levels between the root and the leaves create a severe \textbf{vanishing gradient} problem.
    \item \textbf{CalcParent}: Near-zero gradients (same reason).
    \item \textbf{Emb2Logits (leaf head)}: Receives most of the gradient. But it can only learn the final mapping, not fix upstream representation issues.
\end{itemize}

For BEMAC, this is not an issue because the discrete embedding lookup (\texttt{nn.Embedding}) needs minimal gradient---it only has 3 entries. For GMAC, the continuous MLP front-end needs substantial gradient to learn the mapping from $\mathbb{R}$ to $\mathbb{R}^d$, and it never receives it.

\subsection{Key Discovery: CalcParent Call Counts}

We initially hypothesized that Class B fails because it requires more CalcParent calls (bottom-up operations). Careful counting revealed:
\begin{itemize}[nosep]
    \item $N = 1024$, Class C: 2035 CalcParent calls.
    \item $N = 1024$, Class B: 2033 CalcParent calls.
\end{itemize}
The counts are nearly identical. \textbf{CalcParent frequency is not the bottleneck.} The problem lies elsewhere---likely in the continuous front-end's inability to produce sufficiently precise embeddings.

\subsection{Specific Questions}

\begin{enumerate}
    \item \textbf{Injecting polarization structure}: Is there a way to encode the polarization hierarchy into the neural architecture? For example, can we use different precision levels at different tree depths, or use a loss function that explicitly rewards polarization-like behavior?

    \item \textbf{Loss function redesign}: Cross-entropy treats all bit errors equally. But in SC decoding, an error at position $t$ causes cascade errors at positions $t+1, \ldots, 2N$. Can we design a loss that weights early positions more, or that captures error propagation?

    \item \textbf{Alternative architectures}: Should we abandon the tree-walk structure for continuous channels? A transformer that processes the full observation sequence $Z_1, \ldots, Z_N$ autoregressively might handle continuous inputs better. We tried this for BEMAC Class B at $N = 8$ (matched SC) but it failed at $N = 16$ due to exposure bias.

    \item \textbf{Single-user neural decoders}: The NPD literature (CRISP~\cite{crisp2023}, DeepPolar~\cite{deeppolar2024}) achieves near-ML decoding for single-user polar codes. Can their techniques (e.g., multi-round belief propagation, learned code construction) be extended to the two-user MAC?

    \item \textbf{SCL on top of NN} (\textcolor{green!60!black}{\textbf{RESOLVED}}): We implemented Neural SCL ($L = 4$) on top of the neural tree and it works \emph{spectacularly} for BEMAC: 8.4$\times$ better than SC at $N = 64$. However, for GMAC the neural tree embeddings are still imprecise at large $N$, so the SCL gain is limited. The question shifts to: how can we make the continuous front-end produce sufficiently precise embeddings for the neural SCL to exploit?

    \item \textbf{Is shared-weight the bottleneck?}: All CalcLeft/CalcRight/CalcParent MLPs share weights across tree levels. Does the continuous channel require level-specific parameters? If so, can we achieve this without $\mathcal{O}(n)$ parameter growth (e.g., via hypernetworks or FiLM conditioning on depth)?
\end{enumerate}

\subsection{Available Resources}

\begin{itemize}[nosep]
    \item Mac with 8\,GB RAM (CPU training, $\sim$1 second/iteration at $N = 128$).
    \item Google Colab (T4 GPU, 12-hour free sessions).
    \item Complete analytical SC decoder for GMAC (can generate unlimited training data).
    \item Pre-trained BEMAC NCG decoder that matches/beats SC up to $N = 1024$.
    \item Pre-trained GMAC NCG decoder (works at $N = 32$, fails at $N \ge 64$).
    \item Approximately one week of compute budget.
\end{itemize}

\subsection{Summary}

We have a neural SC decoder architecture that \textbf{provably works} for discrete MAC channels (BEMAC) at all block lengths up to $N = 1024$. Furthermore, \textbf{Neural SCL} ($L = 4$) achieves dramatic gains over both SC and SCL-32 for BEMAC: 8.4$\times$ better than SC at $N = 64$, marking the first successful neural list decoder for MAC polar codes. The same architecture fails for continuous channels (GMAC) at $N \ge 64$. The failure mode is a flat BLER curve that does not benefit from polarization. Gradient analysis points to vanishing gradients through the tree as the proximate cause, and the continuous front-end's inability to produce high-precision embeddings as the root cause. We seek architectural or training innovations that can bridge this gap---particularly for the continuous front-end, so that Neural SCL can deliver similar gains on GMAC as it does on BEMAC.

\begin{thebibliography}{9}
\bibitem{arikan2009}
E.~Ar{\i}kan, ``Channel polarization: A method for constructing capacity-achieving codes for symmetric binary-input memoryless channels,'' \emph{IEEE Trans.\ Inf.\ Theory}, vol.~55, no.~7, pp.~3051--3073, Jul.~2009.

\bibitem{onay2013}
S.~\"Onay, ``Successive cancellation decoding of polar codes for the two-user binary-input MAC,'' in \emph{Proc.\ IEEE Int.\ Symp.\ Inf.\ Theory (ISIT)}, Jul.~2013, pp.~1532--1536.

\bibitem{ren2025}
Y.~Ren \emph{et al.}, ``Polar codes for arbitrary monotone chain rules on MACs,'' \emph{IEEE Trans.\ Inf.\ Theory}, 2025.

\bibitem{crisp2023}
A.~Hebbar \emph{et al.}, ``CRISP: Curriculum based sequential neural decoders for polar codes,'' in \emph{Proc.\ IEEE ICC}, 2023.

\bibitem{deeppolar2024}
S.~Cammerer \emph{et al.}, ``DeepPolar: Neural polar decoders,'' in \emph{Proc.\ IEEE GLOBECOM}, 2024.
\end{thebibliography}

\end{document}
